\section{文献挖掘与信息抽取模块}
\subsection{数据来源与抓取策略}
\placeholder{说明使用 PubMed、关键词设计、抓取字段（PMID、Title、Abstract 等）、脚本文件名。}

\subsection{实体识别与关系抽取策略}
\placeholder{说明抽取的实体类型（TE、Disease、Function、Trait、Paper）和关系类型，说明使用的规则、Prompt 或模型。}

\subsection{数据清洗与规范化}
\placeholder{说明去重、别名统一、LINE-1 亚家族关系补充、JSONL/CSV 到图数据库导入前的处理。}

\subsection{文本挖掘策略与准确率评估（简述）}
\placeholder{这是 PDF 明确要求的内容。建议写人工抽样评估：抽样多少篇摘要、看实体抽取是否正确、关系抽取是否合理。}

\begin{table}[H]
  \centering
  \caption{人工抽样准确率评估模板}
  \begin{tabular}{p{3cm} p{3cm} p{3cm} p{4cm}}
    \toprule
    评估对象 & 样本量 & 正确数 & 备注 \\
    \midrule
    实体识别 & \placeholder{例如 50} & \placeholder{例如 43} & \placeholder{说明判断标准} \\
    关系抽取 & \placeholder{例如 30} & \placeholder{例如 24} & \placeholder{说明判断标准} \\
    \bottomrule
  \end{tabular}
\end{table}

